% TOP_PROBLEM: {"id": 1140, "title": "The Birch and Swinnerton-Dyer Conjecture", "category_id": 5, "category": {"id": 5, "name": "algebraic_geometry", "display_name": "Algebraic Geometry", "description": "Geometric objects defined by polynomial equations.", "slug": "algebraic-geometry", "order_index": 5, "created_at": "2026-07-31T15:26:25.671Z"}, "status": "open"}
% FIRST_POSTED: null
% ENTRY_KIND: statement_only
% Imported from: batch04.tex; UnsolvedMath ID 1140
\documentclass[11pt]{article}
\usepackage[margin=25mm]{geometry}
\usepackage{fontspec}
\setmainfont{Latin Modern Roman}
\usepackage{amsmath,amssymb,mathtools}
\usepackage{unicode-math}
\setmathfont{Latin Modern Math}
\usepackage{xurl,hyperref}
\hypersetup{colorlinks=true,urlcolor=blue}
\setcounter{secnumdepth}{0}
\setlength{\parindent}{0pt}
\setlength{\parskip}{5pt}
\setlength{\emergencystretch}{3em}
\newcommand{\sref}[2]{\hyperlink{source:#1:#2}{[S#2]}}
\newcommand{\cataloguescope}[1]{\paragraph{Recorded target.}#1}
\begin{document}
% UnsolvedMath ID 1140; source catalogue code AG-003
\setcounter{section}{1139}
\section{The Birch and Swinnerton-Dyer Conjecture}\label{1140}
{\small\sffamily UnsolvedMath ID 1140 · Algebraic Geometry}
\subsection{Definitions and mathematical statement}

\paragraph{Shared notation.}$\mathbb N=\{1,2,\ldots\}$, and $\mathbb Z,\mathbb Q,\mathbb R,\mathbb C$ denote the integers, rationals, reals and complex numbers. Any other conventions are specified locally.


An elliptic curve \(E/\mathbb Q\) is a smooth projective genus-one
curve with a specified rational point as identity.  The Mordell--Weil
group has the form
\[
 E(\mathbb Q)\cong E(\mathbb Q)_{\mathrm{tors}}\oplus\mathbb Z^{r_E};
\]
\(r_E\) is its algebraic rank.  Let \(L(E,s)\) be its Hasse--Weil
\(L\)-function.  For a prime \(p\) of good reduction,
\(a_p=p+1-|E(\mathbb F_p)|\), and the local factor is
\((1-a_pp^{-s}+p^{1-2s})^{-1}\).
For split multiplicative, nonsplit multiplicative and additive
reduction, respectively, put \(a_p=1,-1,0\); the local factor is
\((1-a_pp^{-s})^{-1}\).  The product of these factors, initially
convergent for \(\operatorname{Re}(s)>3/2\), has its analytic
continuation.  Its analytic rank is the order of vanishing at \(s=1\):
the integer \(r_{\mathrm{an}}\ge0\) for which
\(L(E,s)=(s-1)^{r_{\mathrm{an}}}u(s)\) with \(u\) holomorphic and
\(u(1)\ne0\).
\paragraph{Statement recorded in the supplied source.}
For every elliptic curve \(E/\mathbb Q\), is
\[
 \operatorname{rank}_{\mathbb Z}E(\mathbb Q)
       =\operatorname{ord}_{s=1}L(E,s)?
\]
This is the rank-equality assertion recorded here; no additional
leading-coefficient formula is included in this question. \sref{1140}{2}
\subsection{Short English statement}
For every rational elliptic curve, does the number of independent rational points equal the multiplicity of the zero of its L-function at one?
\subsection{Sources}
\begin{description}
\item[\hypertarget{source:1140:1}{S1}] ULAM.AI and the UnsolvedMath Contributors, \emph{UnsolvedMath: A Curated Collection of Open Mathematics Problems}, record 1140 (AG-003). CC BY 4.0. \url{https://huggingface.co/datasets/ulamai/UnsolvedMath}.
\item[\hypertarget{source:1140:2}{S2}] Fabio Ferrari Ruffino, \emph{Introduction to the Birch and Swinnerton-Dyer Conjecture}, arXiv:2608.25975 (2026), Sections 4--6, especially the rank conjecture in Section 6. \url{https://arxiv.org/abs/2608.25975}.
\end{description}
\label{end:1140}
\end{document}
